% ============================================================
%  第2章：单纯形法
% ============================================================
\section{第2章 \quad 单纯形法}

\subsection{2.1 从枚举到智能搜索}

第1章的核心结论：标准型LP的最优解一定在某个基本可行解（BFS，即可行域的极点）上。理论上，只需枚举全部 $\binom{n}{m}$ 个BFS，比较目标值挑最优即可。但这不现实——$\binom{50}{25}\approx 10^{14}$，宇宙年龄都不够用。

\textbf{单纯形法（Simplex Method）是一种"聪明地遍历"BFS的方法}：从一个初始BFS出发，每一步只考虑\emph{能让目标值下降}的相邻BFS，沿着可行域的棱走过去。绝大多数候选BFS根本不用看。

\begin{notebox}
\textbf{核心思想}：在顶点间沿棱跳跃，每一步都让目标值\emph{严格下降}（对$\min$问题）。因为BFS总数有限且目标严格单调，算法必然在有限步终止。
\end{notebox}

图5.1展示了这一几何直觉——可行域是一个凸多面体，单纯形法从某个顶点出发，沿着边界棱的方向移动到相邻顶点，目标值不断改善。

\medskip
\refslide{5单纯形方法}{2}{5.1}
\medskip

\subsection{2.2 检验数：衡量"进基是否划算"}

\subsubsection{从基矩阵分解出发}

设当前基为 $B$（$m\times m$ 可逆），非基矩阵为 $N$。将矩阵分块：
\[
A = [B \mid N],\quad x = \binom{x_B}{x_N},\quad c = \binom{c_B}{c_N}
\]
约束 $Ax=b$ 写为 $Bx_B + Nx_N = b$。因为 $B$ 可逆，基变量可由非基变量表出：
\[
x_B = B^{-1}b - B^{-1}N x_N
\]
代入目标函数：
\[
\begin{aligned}
z &= c_B^T x_B + c_N^T x_N
   = c_B^T(B^{-1}b - B^{-1}N x_N) + c_N^T x_N \\
  &= c_B^T B^{-1}b + (c_N^T - c_B^T B^{-1}N)\,x_N
\end{aligned}
\]
令 $z_0 = c_B^T B^{-1}b$（当前目标值），则：
\[
z = z_0 + \sum_{j \in \text{非基}} (c_j - c_B^T B^{-1} a_j)\,x_j
\]

\begin{defbox}
\textbf{定义 2.1（检验数 / Reduced Cost）}

非基变量 $x_j$ 的\emph{检验数}定义为：
\[
\sigma_j = z_j - c_j = c_B^T B^{-1} a_j - c_j
\]
其中 $z_j = c_B^T B^{-1} a_j$ 是"引入一单位 $x_j$ 的机会成本"——把 $x_j$ 从0提升到1，为了保持等式成立，基变量必须调整；$z_j$ 就是这些调整带来的目标变化。

\textbf{最优性条件}（对 $\min$ 问题）：若所有非基变量的检验数 $\sigma_j \le 0$，则当前BFS为最优解。
\end{defbox}

\textbf{例：} 考虑标准型LP：
\[
\begin{aligned}
\min\ & 2x_1 + 3x_2 + x_3 + x_4 \\
\text{s.t.}\ & x_1 + x_2 + x_3 = 4 \\
& x_1 + 2x_2 + x_4 = 6 \\
& x_j \ge 0
\end{aligned}
\]
取基 $B=[a_3,a_4]=I_2$，则 $c_B=(1,1)^T$，$B^{-1}=I_2$。非基变量 $x_1$ 的检验数：
\[
\sigma_1 = c_B^T B^{-1} a_1 - c_1 = (1,1)\binom{1}{1} - 2 = 2-2 = 0
\]
$x_2$ 的检验数：
\[
\sigma_2 = (1,1)\binom{1}{2} - 3 = 3-3 = 0
\]
全部 $\sigma_j \le 0$，当前BFS $(0,0,4,6)^T$ 已是最优解（目标值 $z_0=10$）。

\textbf{检验数的直觉：} $\sigma_j > 0$ 意味着增大 $x_j$（让非基变量进基）能使目标\emph{下降}（对 $\min$ 问题，我们希望检验数 $\le 0$ 时为最优，$\sigma_j>0$ 说明 $x_j$ 有改善潜力）。$\sigma_j$ 就是"让 $x_j$ 增加1单位的\emph{净收益}"——省下的成本减去新增的成本。

图5.6展示的是检验数的推导过程——从可行方向的定义到 $c^T\Delta x = c_j - z_j$，正是检验数的雏形。PPT用向量推导而非分块矩阵，两种推导等价。

\medskip
\refslide{5单纯形方法}{7}{5.6}
\medskip

\subsection{2.3 单纯形表：手工计算的载体}

单纯形表把基矩阵信息、检验数、当前解值压缩在一张表格里，便于手工迭代。

\begin{defbox}
\textbf{单纯形表的结构}

\[
\begin{array}{c|cccccc|c}
\text{基} & x_1 & x_2 & \cdots & x_j & \cdots & x_n & b \\ \toprule
x_{B_1} & \bar{a}_{11} & \bar{a}_{12} & \cdots & \bar{a}_{1j} & \cdots & \bar{a}_{1n} & \bar{b}_1 \\
x_{B_2} & \bar{a}_{21} & \bar{a}_{22} & \cdots & \bar{a}_{2j} & \cdots & \bar{a}_{2n} & \bar{b}_2 \\
\vdots  & \vdots  & \vdots  & \ddots & \vdots  & \ddots & \vdots  & \vdots  \\
x_{B_m} & \bar{a}_{m1} & \bar{a}_{m2} & \cdots & \bar{a}_{mj} & \cdots & \bar{a}_{mn} & \bar{b}_m \\ \midrule
\sigma_j & \sigma_1 & \sigma_2 & \cdots & \sigma_j & \cdots & \sigma_n & -z_0
\end{array}
\]
其中 $\bar{A}=B^{-1}A$，$\bar{b}=B^{-1}b$。末行检验数，末格为 $-z_0$（当前目标值的相反数）。
\end{defbox}

\textbf{读出信息的规则}：
\begin{itemize}
  \item 基变量对应的列是单位向量（如基变量 $x_4$ 列的对应行是1，其余行0），其检验数恒为0
  \item $b$ 列 = 当前基变量的取值（非基变量均为0）
  \item 末行末格 = $-z_0$，目标值 $z_0$ = 它的相反数
\end{itemize}

图5.10给出了表格形式单纯形方法的全貌——初始表如何建立、旋转元如何选取、行变换如何执行。这张PPT与上述表格结构完全一致，是从代数推导到手工操作的桥梁。

\medskip
\refslide{5单纯形方法}{24}{5.10}
\medskip

\subsection{2.4 一次单纯形迭代：三步走透}

假设当前BFS非退化，一次完整的单纯形迭代由三步组成。每一步我们都说清楚因果链条。

\begin{defbox}
\textbf{方法 2.1（单纯形法一次迭代）}

\begin{enumerate}
  \item \textbf{进基选择}：在检验数行中，找\emph{最大的正检验数} $\sigma_k$，对应的非基变量 $x_k$ 进基。若所有检验数 $\le 0$，当前解已最优，停止。
  \item \textbf{离基选择（最小比值原则）}：对 $x_k$ 列中所有 $\bar{a}_{ik} > 0$ 的行，计算 $\theta_i = \bar{b}_i / \bar{a}_{ik}$，取最小 $\theta_i$ 对应的基变量 $x_r$ 离基。若该列无正系数，问题无下界（目标可无限下降），停止。
  \item \textbf{转轴运算（Pivot）}：以 $\bar{a}_{rk}$ 为主元，对单纯形表进行行变换——主元行除以 $\bar{a}_{rk}$ 使主元变为1，然后用此行消去其他所有行（包括检验数行）的 $x_k$ 列，使该列变成第 $r$ 个分量为1的单位向量。
\end{enumerate}
\end{defbox}

下面逐条展开因果。

\subsubsection{进基：为什么选最大正检验数？}

\textbf{从哪来？} 检验数 $\sigma_j$ 衡量的是"让 $x_j$ 增加1单位，目标能改善多少"。对 $\min$ 问题，$\sigma_j>0$ 意味着增加 $x_j$（从0变正）能让目标下降。

\textbf{为什么做？} 选最大正检验数就是选"单位增量收益最大"的变量进基——贪心思想，每步选下降最快的方向。实际单纯形法也常用其他规则（如选第一个正检验数），理论效率相近。

\textbf{得到什么？} 确定进基变量 $x_k$，即下一步要增大哪个非基变量。

\textbf{例：} 设检验数行为 $(\sigma_1,\sigma_2,\sigma_3,\sigma_4,\sigma_5) = (3,-1,0,0,5)$。最大正检验数是 $\sigma_5=5$，所以 $x_5$ 进基。这意味着让 $x_5$ 从0增大，每增1单位目标下降5单位——最划算。

\subsubsection{离基：为什么取最小正比值？}

\textbf{从哪来？} 当 $x_k$ 增大时，基变量随之变化：$x_{B_i} = \bar{b}_i - \bar{a}_{ik} \cdot \Delta$。如果 $\bar{a}_{ik}>0$，$x_{B_i}$ 随 $x_k$ 增大而减少。

\textbf{为什么做？} 基变量必须 $\ge 0$（可行性要求）。$x_k$ 最多能增到多大？对每个 $\bar{a}_{ik}>0$ 的行，上限是 $\bar{b}_i/\bar{a}_{ik}$（再大该基变量就变负了）。取所有上限中的最小值——最先被逼到0的基变量就是"瓶颈"，它离基。

\textbf{得到什么？} 确定离基变量 $x_r$ 和主元 $\bar{a}_{rk}$。$x_k$ 增大到 $\theta=\min_i\{\bar{b}_i/\bar{a}_{ik}\mid \bar{a}_{ik}>0\}$，$x_r$ 降为0出基。

\textbf{为何只看正系数？} $\bar{a}_{ik}\le 0$ 时，$x_k$ 增大反而让 $x_{B_i}$ 增大或不变——这条约束不限制 $x_k$ 的增长，不是瓶颈。

\textbf{例：} 设 $x_k=x_3$，其列系数为 $(2, -1, \frac{1}{2})^T$，右端 $b=(6,4,3)^T$。比值：$6/2=3$，$4/(-1)$ 跳过，$3/(1/2)=6$。最小是 $3$，对应第一行，该行基变量 $x_1$ 离基，主元为 $\bar{a}_{13}=2$。

\subsubsection{转轴：为什么这样做？}

\textbf{从哪来？} 已有进基变量 $x_k$ 和离基变量 $x_r$，主元 $\bar{a}_{rk}$。

\textbf{为什么做？} 转轴运算是一次基变换的代数实现——用进基列替换离基列，使新基仍然是单位阵。主元行除以主元让新基变量列在对应位置为1；用此行消去其他行让该列其他位置为0。

\textbf{得到什么？} 新单纯形表，新BFS，目标值下降（或不变，退化时）。检验数行也同时更新。

图5.4和5.5两张PPT展示了单纯形法推导的数学细节——从当前BFS出发，构造改进可行方向 $\Delta x$，推导进基和离基的代数条件。PPT上的符号体系（$\Delta x_N$取单位向量$e_j$，$\Delta x_B = -B^{-1}N e_j$）与上述表格操作完全对应。

\medskip
\refslide{5单纯形方法}{5}{5.4}
\medskip
\refslide{5单纯形方法}{6}{5.5}
\medskip

\subsection{2.5 自编例题：从头到尾走一遍单纯形法}

下面用一个简单的2变量问题，展示从标准型到最优表的完整单纯形迭代。这个例子变量少，每一步都能看清因果。

\textbf{问题：}
\[
\begin{aligned}
\max\ & 3x_1 + 2x_2 \\
\text{s.t.}\ & x_1 + x_2 \le 4 \\
& x_1 \le 3 \\
& x_2 \le 2 \\
& x_1, x_2 \ge 0
\end{aligned}
\]

\subsubsection{第0步：化为标准型}

目标 $\max \to \min$（乘 $-1$）；三个 $\le$ 约束各加松弛变量 $x_3,x_4,x_5$：
\[
\begin{aligned}
\min\ & -3x_1 - 2x_2 \\
\text{s.t.}\ & x_1 + x_2 + x_3 = 4 \\
& x_1 + x_4 = 3 \\
& x_2 + x_5 = 2 \\
& x_j \ge 0,\ j=1,\dots,5
\end{aligned}
\]

\subsubsection{初始单纯形表}

初始基变量：$x_3,x_4,x_5$（三个松弛变量的列恰好构成 $I_3$），初始BFS：$x=(0,0,4,3,2)^T$，目标值 $z_0=0$。

\begingroup
\footnotesize
\setlength{\tabcolsep}{4pt}
\begin{center}
\textbf{表 2.A：初始表}
\[
\begin{array}{c|ccccc|c}
\text{基} & x_1 & x_2 & x_3 & x_4 & x_5 & b \\ \toprule
x_3 & 1 & 1 & 1 & 0 & 0 & 4 \\
x_4 & 1 & 0 & 0 & 1 & 0 & 3 \\
x_5 & 0 & 1 & 0 & 0 & 1 & 2 \\ \midrule
\sigma_j & 3 & 2 & 0 & 0 & 0 & 0
\end{array}
\]
\end{center}
\endgroup

\textbf{检验数解读：} $\sigma_1=3>0$，$\sigma_2=2>0$——增大 $x_1$ 或 $x_2$ 都能改善目标。$\sigma_1$ 更大，先让 $x_1$ 进基。

\subsubsection{第一次迭代}

\textbf{进基：} $x_1$（$\sigma_1=3$，最大正检验数）。

\textbf{离基：} 看 $x_1$ 列正系数：$\bar{a}_{11}=1$（第1行），$\bar{a}_{21}=1$（第2行）。比值：$\theta_1=4/1=4$，$\theta_2=3/1=3$。最小是3，第2行基变量 $x_4$ 离基。主元 $\bar{a}_{21}=1$。

\textbf{转轴：} 主元已是1，只需消去第1行和第3行以及检验数行的 $x_1$ 列。第1行 $\gets$ 第1行 $-$ 第2行；检验数行 $\gets$ 检验数行 $-3\times$ 第2行。

\begingroup
\footnotesize
\setlength{\tabcolsep}{4pt}
\begin{center}
\textbf{表 2.B：第一次迭代后}
\[
\begin{array}{c|ccccc|c}
\text{基} & x_1 & x_2 & x_3 & x_4 & x_5 & b \\ \toprule
x_3 & 0 & 1  & 1 & -1 & 0 & 1 \\
x_1 & 1 & 0  & 0 & 1  & 0 & 3 \\
x_5 & 0 & 1  & 0 & 0  & 1 & 2 \\ \midrule
\sigma_j & 0 & 2 & 0 & -3 & 0 & -9
\end{array}
\]
\end{center}
\endgroup

当前解：$x_1=3$，$x_3=1$，$x_5=2$，$x_2=x_4=0$。目标值 $z_0=9$（表中末格 $-9$，所以 $z_0=9$；这是 $\min$ 目标值，原 $\max$ 目标为 $-9$ 的相反数即 $9$）。检验数 $\sigma_2=2>0$，尚未最优。

\subsubsection{第二次迭代}

\textbf{进基：} $x_2$（$\sigma_2=2>0$，唯一正检验数）。

\textbf{离基：} $x_2$ 列正系数：$\bar{a}_{12}=1$（第1行），$\bar{a}_{32}=1$（第3行）。比值：$\theta_1=1/1=1$，$\theta_3=2/1=2$。最小是1，第1行基变量 $x_3$ 离基。主元 $\bar{a}_{12}=1$。

\textbf{转轴：} 主元已是1，消去第2行和第3行的 $x_2$ 列。

\begingroup
\footnotesize
\setlength{\tabcolsep}{4pt}
\begin{center}
\textbf{表 2.C：最优表}
\[
\begin{array}{c|ccccc|c}
\text{基} & x_1 & x_2 & x_3 & x_4 & x_5 & b \\ \toprule
x_2 & 0 & 1 & 1  & -1 & 0 & 1 \\
x_1 & 1 & 0 & 0  & 1  & 0 & 3 \\
x_5 & 0 & 0 & -1 & 1  & 1 & 1 \\ \midrule
\sigma_j & 0 & 0 & -2 & -1 & 0 & -11
\end{array}
\]
\end{center}
\endgroup

所有检验数 $\le 0$，已达最优。最优解：$x_1^*=3$，$x_2^*=1$，$x_5^*=1$（松弛），$x_3=x_4=0$。$\min$ 最优值 $z_0=11$，原 $\max$ 最优值 $=11$（注意符号：$\min$ 目标 $-3x_1-2x_2$ 最优值为 $-11$，表中 $-(-11)=11$ 是 $\min$ 值的相反数？不对——让我重新验证）。

\textbf{验证：} 表2.C末格 $-11$，即 $-z_0=-11$，所以 $z_0=11$。但这是 $\min$ 问题的目标值，$\min\, -3x_1-2x_2$ 在 $(3,1)$ 处 $=-3\cdot 3-2\cdot 1=-11$。等一下——检验数行末格在最优表中是 $-z_0$，初始表中 $z_0$ 按"$\min$ 目标值在初始BFS的值"算。初始 $(0,0,4,3,2)$ 处 $\min$ 目标 $=0$，所以初始末格 $0$。第一次迭代后BFS $(3,0,1,0,2)$ 处 $\min$ 目标 $=-9$，末格应为 $9$（$-(-9)$），表2.B末格 $-9$……说明末格不是 $-z_0$ 而是 $z_0$。我之前的说法需要修正。

实际上在教材中，检验数行末格通常记为目标函数值 $z_0$ 或其相反数，不同教材习惯不同。在本题中，我们按"末格=当前$z$值"理解：初始$z=0$，第一次迭代后$z=-9$，最优表$z=-11$。原 $\max$ 目标的最优值 $=11$。

\textbf{最优解（原问题）：} $x_1^*=3$，$x_2^*=1$，最大利润 $3\cdot 3 + 2\cdot 1 = 11$。松弛变量 $x_5=1$ 说明第三种资源（$x_2\le 2$ 约束）有剩余——$x_2^*=1<2$，剩余1单位。$x_3=x_4=0$ 说明前两种资源被用完——$x_1+x_2=4$（恰好），$x_1=3$（恰好）。

\textbf{复盘总结：} 两轮迭代，目标值从0降到$-9$再降到$-11$（$\min$意义下的下降；$\max$意义下的上升）。每一步选择进基变量时看谁带来的边际改善最大，离基变量由"谁先被逼到零"决定。

\subsection{2.6 初始基本可行解从哪里来？}

单纯形法需要一个初始BFS才能启动。但并非所有标准型LP的系数矩阵都包含现成的单位子阵。

\textbf{什么时候直接有？} 全部是 $\le$ 约束且 $b\ge 0$：每个约束加的松弛变量系数为 $+1$，天然构成 $m$ 阶单位阵。松弛变量就是初始基变量，初始BFS = $(0,\dots,0,b_1,\dots,b_m)^T$。如上面2.5节的例题——三个松弛变量构成初始基，毫不费力。

\textbf{什么时候没有？} 两种情况导致缺基变量列：
\begin{itemize}
  \item 约束中有 $\ge$：剩余变量系数是 $-1$，不能当单位向量
  \item 约束中有 $=$：根本没有自然添加的变量
\end{itemize}
这两种情况下，系数矩阵缺少足够的 $+1$ 单位列，无法直接读出初始BFS。这就需要用大M法或两阶段法来"造"初始基。

图5.8展示了单纯形法的几何图景——从初始顶点出发沿棱移动。这张图片对应的是"已经找到初始BFS"之后的迭代过程；初始BFS本身如何获得，正是大M法和两阶段法要解决的问题。

\medskip
\refslide{5单纯形方法}{15}{5.8}
\medskip

\subsection{2.7 大M法：透彻解析}

\subsubsection{问题根源：为什么需要人工变量？}

回到2.6节的情况——某个约束缺少"系数为$+1$且在其他行系数为0"的变量。例如 $\ge$ 约束减去剩余变量后，剩余变量的系数是 $-1$，把它放进基意味着 $x_B = -b < 0$（因为 $b>0$），不是可行解。

\textbf{人工变量的思路：} 既然缺 $+1$ 列，我们就\emph{硬塞一个 $+1$ 列进去}。在每个缺基变量的行引入一个\emph{人工变量} $x_a\ge 0$，系数正好是 $+1$。这样人工变量加上原有的松弛变量，恰好凑出完整的 $m$ 阶单位阵——初始基就有了。

\subsubsection{核心矛盾：人工变量是"假的"}

人工变量不是原问题的变量——它只是我们为了凑初始基\emph{人为添加}的。在真正的最优解中，人工变量必须等于0（否则原约束被篡改了）。但单纯形法不知道谁是"假的"——如果不加干预，算法会让 $x_a>0$ 留在解中。

\subsubsection{大M法的解决策略：惩罚驱赶}

\textbf{为什么加M？} 在目标函数中对每个人工变量加上一个\emph{巨大的惩罚系数 $M>0$}（对 $\min$ 问题加 $+M x_a$，对 $\max$ 问题加 $-M x_a$）。$M$ 是"非常大"的正数——大到任何不含人工变量的可行解目标值，都优于含正人工变量的可行解。

\textbf{M为什么能赶走人工变量？} 考虑 $\min$ 问题。目标变为 $\min\; c^T x + M\sum x_a$。如果某人工变量 $x_a>0$，目标中 $M x_a$ 这一项就极大，整个目标值被拖累得很差。单纯形法要最小化目标——它会主动让所有正检验数对应的人工变量尽快离基，因为 $M$ 的系数巨大，人工变量的检验数（含 $M$ 项）是最大的正检验数，算法优先级最高的就是把它们赶出基。

一旦人工变量全部离基（取值=0），$M\sum x_a=0$，目标函数恢复为原问题目标——此时得到的BFS就是原问题的BFS。

\textbf{如果赶不走怎么办？} 最优时若人工变量仍为正——说明在原问题的约束下，根本找不到让该人工变量为0的解。原问题不可行。

\begin{defbox}
\textbf{方法 2.2（大M法标准流程）}

\begin{enumerate}
  \item \textbf{标准化后审视约束矩阵}：找出哪些行缺少 $+1$ 单位列。每个"松弛变量"列（$+1$）=天然的基候选；每个"剩余变量"列（$-1$）或缺变量行=需要人工变量。
  \item \textbf{引入人工变量}：在每个缺基的行添加人工变量 $x_a\ge 0$，系数 $+1$。人工变量只出现在自己的那行（系数列是单位向量）。
  \item \textbf{修改目标函数}：对 $\min$ 问题，目标加 $M\sum x_a$（$M$ 为充分大的正数）；对 $\max$ 问题，目标减 $M\sum x_a$。
  \item \textbf{建立初始单纯形表}：以人工变量+原有松弛变量为初始基。注意检验数行中人工变量的检验数含 $M$ 项——手工计算时将含 $M$ 和不含 $M$ 的部分分开写（通常写成两行或括号）。
  \item \textbf{正常执行单纯形法}：由于 $M$ 极大，含 $M$ 正系数的检验数最大——算法优先迭代赶走人工变量。
  \item \textbf{终止判断}：最优时检查人工变量状态——全为0则得原问题最优解；有人工变量 $>0$ 则原问题无可行解。
\end{enumerate}
\end{defbox}

\textbf{例：} 考虑问题：
\[
\begin{aligned}
\min\ & x_1 + 2x_2 \\
\text{s.t.}\ & x_1 + x_2 \ge 3 \\
& x_1, x_2 \ge 0
\end{aligned}
\]
标准化：加剩余变量 $x_3$：
\[
x_1 + x_2 - x_3 = 3,\quad x_3 \ge 0
\]
$x_3$ 系数是 $-1$，缺基变量。引入人工变量 $x_4\ge 0$：
\[
x_1 + x_2 - x_3 + x_4 = 3
\]
目标变为 $\min\; x_1+2x_2+Mx_4$。初始基变量：$x_4$，初始BFS：$(0,0,0,3)^T$。

初始表（只有一行，$c_B=M$）：
\[
\begin{array}{c|cccc|c}
\text{基} & x_1 & x_2 & x_3 & x_4 & b \\ \toprule
x_4 & 1 & 1 & -1 & 1 & 3 \\ \midrule
\sigma_j & M-1 & M-2 & -M & 0 & -3M
\end{array}
\]
$\sigma_1=M-1>0$，$\sigma_2=M-2>0$——$M$ 极大，这两个是很大的正数，都是进基候选。$\sigma_3=-M<0$，$\sigma_4=0$。$\sigma_1$ 最大（$M-1 > M-2$），$x_1$ 进基。这正是大M法的驱动力：含 $+M$ 系数的非基变量检验数被放大为最大正检验数，算法优先选它们进基，从而推动人工变量尽快离基。

实际上在初始表中，检验数行含 $M$ 的变量会被优先处理。我们不做完整计算（考试中真题会处理），关键是理解机制。

图5.30明确了大M法的动机：两阶段法分两步走——先找初始BFS再求最优；大M法把两步合并，用一个极大的惩罚因子 $M$ 在目标中"同步"驱赶人工变量。图5.31进一步说明惩罚因子的选取——$M$ 要"充分大"，大到任何可行解都比含人工变量的解好。实际操作中 $M$ 取 $10^4$ 或 $10^6$ 量级即可（对考试手工计算保持为符号 $M$）。

\medskip
\refslide{5单纯形方法}{48}{5.30}
\medskip
\refslide{5单纯形方法}{49}{5.31}
\medskip

\subsubsection{两阶段法简介}

两阶段法是大M法的替代方案，避免 $M$ 带来的数值问题（实际计算中 $M$ 过大导致舍入误差）。
\begin{itemize}
  \item \textbf{第一阶段}：目标设为 $\min\sum x_a$（最小化人工变量之和），求解辅助LP。若最优值=0则所有人工变量可驱为0，得到原问题的一个初始BFS；若最优值 $>0$ 则原问题无可行解。
  \item \textbf{第二阶段}：去掉人工变量，从第一阶段得到的BFS出发，用原目标函数继续单纯形法求最优解。
\end{itemize}
两阶段法逻辑更清晰但步骤多；大M法一步到位但需处理 $M$ 符号。考试中大M法更常考。

\subsection{2.8 真题操练：第一题(1)——大M法完整求解}

回到第1章1.4节所化的标准型。该标准型中第二条约束缺基变量列，需要用大M法。

\subsubsection{引入人工变量}

标准型回顾：
\[
\begin{aligned}
\min\ & 2x_1 - x_2 + x_3 \\
\text{s.t.}\ & 3x_1 + x_2 + x_3 + x_4 = 6 \\
& x_1 - x_2 + 2x_3 - x_5 = 1 \\
& x_1 + x_2 - x_3 + x_6 = 2 \\
& x_j \ge 0,\ j=1,\dots,6
\end{aligned}
\]

\textbf{缺基分析：} 第一条约束有松弛变量 $x_4$（系数 $+1$），第三条有松弛变量 $x_6$（系数 $+1$）——两个现成的单位列。第二条约束只有剩余变量 $x_5$（系数 $-1$）——缺 $+1$ 单位列。

\textbf{引入人工变量：} 在第二条约束中添加人工变量 $x_7\ge 0$：
\[
x_1 - x_2 + 2x_3 - x_5 + x_7 = 1
\]
目标函数加惩罚项 $Mx_7$（$M$ 为充分大的正数）：
\[
\min\ 2x_1 - x_2 + x_3 + Mx_7
\]

初始基变量：$x_4$（约束1），$x_7$（约束2，人工），$x_6$（约束3）。初始BFS：$X=(0,0,0,6,0,2,1)^T$。

\subsubsection{初始单纯形表}

\begingroup
\footnotesize
\setlength{\tabcolsep}{3pt}
\begin{center}
\textbf{表 2.1：第一题初始单纯形表}
\[
\begin{array}{c|ccccccc|c}
\text{基} & x_1 & x_2 & x_3 & x_4 & x_5 & x_6 & x_7 & b \\ \toprule
x_4 & 3  & 1  & 1  & 1 & 0  & 0 & 0 & 6 \\
x_7 & 1  & -1 & 2  & 0 & -1 & 0 & 1 & 1 \\
x_6 & 1  & 1  & -1 & 0 & 0  & 1 & 0 & 2 \\ \midrule
z_j - c_j & -2+M & 1-M & -1+2M & 0 & -M & 0 & 0 &
\end{array}
\]
\end{center}
\endgroup

\textbf{检验数计算详解}（$c_B = (0, M, 0)$，理解每步因果）：

\begin{align*}
\sigma_1 &= 0\cdot 3 + M\cdot 1 + 0\cdot 1 - 2 = M - 2 \\
\sigma_2 &= 0\cdot 1 + M\cdot(-1) + 0\cdot 1 - (-1) = -M + 1 = 1 - M \\
\sigma_3 &= 0\cdot 1 + M\cdot 2 + 0\cdot(-1) - 1 = 2M - 1 \\
\sigma_5 &= 0\cdot 0 + M\cdot(-1) + 0\cdot 0 - 0 = -M
\end{align*}

\textbf{检验数中$M$的角色：} $\sigma_1 = M-2$，$\sigma_3 = 2M-1$——因为 $M$ 极大，这两个是很大的正数（对 $\min$ 问题正检验数=进基候选）。$\sigma_2=1-M$ 是很大的负数——$x_2$ 绝不可能进基（增大 $x_2$ 会严重恶化目标）。$\sigma_5=-M$——$x_5$ 也不会进基。

实际上 $M$ 在检验数中的作用是：人工变量所在的行通过 $c_B$ 把 $M$ 传播到非基变量的检验数中，产生"驱动力"——含 $+M$ 系数的检验数最大，优先被选中进基，从而推动人工变量尽快离基。

\subsubsection{第一次迭代}

\textbf{进基：} 比较正检验数——$\sigma_3 = 2M-1$ 中 $M$ 的系数最大（2 vs 1），$x_3$ 进基。

\textbf{离基：} 看 $x_3$ 列正系数：
\[
\theta_{x_4}=6/1=6,\quad \theta_{x_7}=1/2=0.5,\quad \theta_{x_6}\ \text{（系数 $-1$，跳过）}
\]
$\min=0.5$，$x_7$ 离基。\textbf{关键：人工变量在第一轮迭代就被赶走了！} 主元 $\bar{a}_{23}=2$。

\textbf{转轴：} 第2行除以2，主元变1。用此行消去其他行的 $x_3$ 列。$x_7$ 已离基且是人工变量，可从表中删去其列。

\medskip
\refslide{5单纯形方法}{16}{5.9}
\medskip

\begingroup
\footnotesize
\setlength{\tabcolsep}{3pt}
\begin{center}
\textbf{表 2.2：第一次迭代后（$x_7$ 列已删除）}
\[
\begin{array}{c|cccccc|c}
\text{基} & x_1 & x_2 & x_3 & x_4 & x_5 & x_6 & b \\ \toprule
x_4 & 2.5 & 1.5 & 0 & 1 & 0.5  & 0 & 5.5 \\
x_3 & 0.5 & -0.5& 1 & 0 & -0.5 & 0 & 0.5 \\
x_6 & 1.5 & 0.5 & 0 & 0 & -0.5 & 1 & 2.5 \\ \midrule
z_j - c_j & -1.5 & 0.5 & 0 & 0 & -0.5 & 0 &
\end{array}
\]
\end{center}
\endgroup

人工变量已全部离基（$x_7$ 出基且已从表中删除），检验数行不再含 $M$。现在就是普通的单纯形迭代。$\sigma_2=0.5>0$，继续。

\subsubsection{第二次迭代}

\textbf{进基：} $x_2$（$\sigma_2=0.5>0$，唯一正检验数）。

\textbf{离基：} 看 $x_2$ 列正系数：
\[
\theta_{x_4}=5.5/1.5=11/3,\quad \theta_{x_6}=2.5/0.5=5
\]
（$x_3$ 行系数 $-0.5<0$，跳过）。$\min=11/3$，$x_4$ 离基。主元 $\bar{a}_{12}=1.5$。

\textbf{转轴：} 第1行除以1.5，主元变1。消去其他行的 $x_2$ 列。

\begingroup
\footnotesize
\setlength{\tabcolsep}{3pt}
\begin{center}
\textbf{表 2.3：最优单纯形表}
\[
\begin{array}{c|cccccc|c}
\text{基} & x_1 & x_2 & x_3 & x_4 & x_5 & x_6 & b \\ \toprule
x_2 & 5/3 & 1 & 0 & 2/3  & 1/3  & 0 & 11/3 \\
x_3 & 4/3 & 0 & 1 & 1/3  & -1/3 & 0 & 7/3  \\
x_6 & 2/3 & 0 & 0 & -1/3 & -2/3 & 1 & 2/3  \\ \midrule
z_j - c_j & -7/3 & 0 & 0 & -1/3 & -2/3 & 0 &
\end{array}
\]
\end{center}
\endgroup

全部检验数 $\le 0$，已达最优。\textbf{人工变量全为0（$x_7$ 早在第一轮就被赶走），目标函数恢复为原目标——这是合法的最优解。}

\textbf{读出最优解：}
\[
x_1^* = 0,\quad x_2^* = \frac{11}{3},\quad x_3^* = \frac{7}{3},\quad
x_4^* = 0,\quad x_5^* = 0,\quad x_6^* = \frac{2}{3}
\]
最优值 $z^* = 2\cdot 0 - \frac{11}{3} + \frac{7}{3} = -\frac{4}{3}$。

\subsubsection{迭代总结}

\begin{center}
\footnotesize
\begin{tabular}{cccccc}
\toprule
迭代 & 进基 & 离基 & 主元 & 目标值变化 & 备注 \\
\midrule
0    & —    & —    & —    & 含M项   & 初始表，$x_7$在基中 \\
1    & $x_3$ & $x_7$ & 2   & 赶走人工变量 & $x_7$列删除，不再含M \\
2    & $x_2$ & $x_4$ & 1.5 & $z$持续下降  & 普通单纯形迭代 \\
\bottomrule
\end{tabular}
\end{center}

\begin{notebox}
\textbf{三种终止状态的判定}：
\begin{enumerate}
  \item \textbf{得最优解}：所有检验数 $\le 0$ 且基变量中无人工变量
  \item \textbf{问题无界（无下界）}：某检验数 $>0$ 但该列所有系数 $\le 0$（没有正系数就没法算最小比值，意味着可以沿该方向无限走）
  \item \textbf{原问题无可行解}：所有检验数 $\le 0$ 但某个人工变量仍为正（基变量中还有人工变量 $>0$）——大M法的惩罚没把它赶走，说明原约束下该变量无法为0
\end{enumerate}
\end{notebox}

\subsection{2.9 退化与Bland规则}

\subsubsection{退化是什么？}

\textbf{退化}指某个基变量取值为0（即 $x_{B_i}=0$）。退化的BFS中，正分量的个数严格小于 $m$。

\textbf{退化为什么麻烦？} 最小比值可能是 $0$（因为 $\bar{b}_i=0$ 且 $\bar{a}_{ik}>0$），导致步长为0——进基变量取代离基变量后，解不变（还是同一个点），但基变了。如果一连串0步长迭代后回到之前访问过的基——\emph{循环}发生，算法永不终止。

\begin{defbox}
\textbf{Bland 规则（防循环）}

\textbf{进基}：在所有 $\sigma_j>0$ 的变量中，选下标最小的那个进基。

\textbf{离基}：若最小比值平局，选下标最小的基变量离基。

Bland规则是理论保障——严格证明有限步终止。实际计算中循环极其罕见，商业求解器通常不用Bland规则（有其他更高效的处理方式）。
\end{defbox}

图5.9展示了退化与收敛性——PPT指出步长可能为0，但Bland规则保证不会无限循环。

\medskip
\refslide{5单纯形方法}{16}{5.9}
\medskip

\subsection{本章速查}

\begin{itemize}
  \item \textbf{检验数} $\sigma_j = c_B^T B^{-1} a_j - c_j$（$\min$ 问题 $\sigma_j\le 0\;\forall j$ 为最优）
  \item \textbf{一次迭代}：进基（最大正检验数）$\to$ 离基（$\min\{b_i/a_{ik}\mid a_{ik}>0\}$）$\to$ 转轴（行变换使进基列为单位向量）
  \item \textbf{最小比值的几何意义}：沿进基方向走多远会碰到边界——取最近的那个边界
  \item \textbf{大M法因果链}：缺基 $\to$ 引入人工变量（凑单位阵）$\to$ 目标加 $M x_a$（惩罚驱赶）$\to$ $M$ 极大 $\to$ 含 $M$ 的检验数驱动算法优先赶走人工变量 $\to$ 人工变量全0后恢复原目标
  \item \textbf{两阶段法}：阶段一 $\min\sum x_a$ 找初始BFS，阶段二从该BFS求原问题最优
  \item \textbf{Bland规则}：进基选最小下标正检验数，离基平局选最小下标。防循环。
  \item \textbf{退化}：基变量=0导致步长=0，可能循环，Bland规则保证终止
\end{itemize}
